%% ============================================================
\section{Related Work}
\label{sec:related}
%% ============================================================

\paragraph{Autonomous research agents.}
End-to-end autonomous research systems have rapidly expanded from constrained ML templates to multi-stage pipelines that coordinate literature grounding, hypothesis generation, experimentation, and paper writing.
The AI Scientist~\citep{lu2024aiscientist} pioneered end-to-end automation but operates on fixed ML templates with frequent hallucinations in writing and limited paper quality.
AI Scientist-v2~\citep{yamada2025aiscientistv2} advances this with best-first tree search (BFTS) over experimental branches and review-aware reporting, achieving workshop-level paper quality.
Concurrent systems extend the pipeline in different directions.
On the ideation side, PiFlow~\citep{wu2025piflow} steers hypothesis exploration via information-theoretic principle selection and CodeScientist~\citep{jansen2025codescientist} grounds ideation jointly in literature and code.
Curie~\citep{kon2025curie} validates experimental execution through reproducibility checks analogous to our I1 Score Verification, though it does not audit whether written claims faithfully reflect the validated results.
Agent Laboratory~\citep{schmidgall2025agentlab} introduces human gating into the pipeline. AlphaEvolve~\citep{novikov2025alphaevolve} applies evolutionary search to algorithmic optimization, and EvoScientist~\citep{li2026evoscientist} uses multi-agent self-evolution for end-to-end discovery.
We evaluate AI Scientist-v2 alongside three additional systems---AutoResearchClaw~\citep{li2024autoresearchclaw}, DeepScientist~\citep{huang2025deepscientist}, and AI-Researcher~\citep{tang2025airesearcher}---whose architectural choices produce distinct integrity profiles (\S\ref{sec:exp:integrity}).
Despite this architectural diversity, a common pattern emerges: generation and execution capabilities have scaled faster than validation and provenance mechanisms, so systems that produce professional-looking manuscripts may still contain broken evidence chains.
\sys{} targets this gap---rather than advancing the autonomy frontier, we focus on making autonomous research outputs verifiable.


\paragraph{LLM-driven optimization and benchmarks.}
The ADRS benchmark~\citep{cheng2025barbarians} collects real frontier computer system research questions and serves as our primary evaluation testbed.
EvoX~\citep{liu2026evox} and AdaEvolve~\citep{cemri2026adaevolve} achieve strong results on ADRS by focusing on algorithm discovery and implementation optimization without literature grounding or paper writing.
Broader evaluation resources have recently proliferated. Auto-Bench~\citep{zhong2025autobench}, ResearchBench~\citep{liu2025researchbench}, and ResearcherBench~\citep{xu2025researcherbench} evaluate research-adjacent capabilities such as causal reasoning, hypothesis generation, and research question answering. MLAgentBench~\citep{huang2024mlagentbench}, EXP-Bench~\citep{weng2025expbench}, and PaperBench~\citep{starace2025paperbench} stress-test experimentation, replication, and execution reliability. AIRS-Bench~\citep{jansen2025airsbench} tests agent performance on tasks drawn from published ML papers. FIRE-Bench~\citep{shojaee2025firebench} evaluates whether agents can rediscover established findings through full-cycle experimentation.
However, most benchmarks measure discovery performance---whether a system can produce competitive solutions---rather than whether the resulting claims are actually supported by evidence.


\paragraph{Scientific integrity and provenance.}
Current autonomous research systems produce written outputs with varying degrees of traceability: direct manuscript drafting where an LLM generates prose from agent outputs~\citep{lu2024aiscientist,jansen2025codescientist,tang2025airesearcher}, and review-aware revision where reviewer feedback refines the manuscript~\citep{yamada2025aiscientistv2}.
Both approaches produce fluent papers but lack mechanisms to ensure that reported numbers trace to specific execution artifacts, masking broken evidence chains.
Prior work on citation verifiability~\citep{liu2023evaluating}, factual accuracy~\citep{min2023factscore}, and citation attribution~\citep{press2024citeme} performs post-hoc detection at the text level.
CoE differs in two ways: it defines verifiability at the level of individual claims (each must trace to a grounding source through the full research artifact), and it covers paper, code, and evaluator logs jointly, not just text.
\coeaudit{} operationalizes this standard as a cross-system audit, subject to the artifact requirements detailed in \S\ref{sec:coeaudit}.
